\clearpage
\onecolumn
\section{Mathematical Definitions and Metric Derivations}
\label{sec:supp-derivations}

This appendix expands the metric definitions used in the main paper. The
definitions are included to make clear which quantities are image-level proxy
diagnostics, which are VLM grounding scores, and which are navigation metrics.

\subsection{Paired Render Notation}

For an asset or scene instance $i$, let $I_i^{o}$ denote the original MDL
render and $I_i^{c}$ denote the converted noMDL render under the matched camera,
resolution, renderer, and lighting configuration. The paired design is
important: all proxy comparisons are computed on pairs whose scene identity and
viewpoint are fixed before conversion. A proxy metric $m$ is reported as

\begin{equation}
\bar{m} = \frac{1}{N}\sum_{i=1}^{N} m(I_i^{o}, I_i^{c}),
\end{equation}

where $N$ is the number of paired render comparisons in that check. This average
is descriptive for the sampled render set; unseen assets require separate
checks.

\begin{center}
\includegraphics[width=\textwidth,height=0.40\textheight,keepaspectratio]{figures/fig_supplement_proxy_metric_derivation_companion.png}
{\captionsetup{hypcap=false,font=scriptsize,skip=2pt}
\captionof{figure}{Proxy metric derivation companion. Paired renders and the
proxy render ladder anchor the proxy equations. The diagram summarizes how the
proxy quantities relate to the sampled render pairs.}}
\label{fig:supp-proxy-metric-derivation-companion}
\end{center}

\subsection{PSNR}

For images normalized to $[0,1]$, mean squared error is

\begin{equation}
\operatorname{MSE}(I^o, I^c) =
\frac{1}{HWC}\sum_{y=1}^{H}\sum_{x=1}^{W}\sum_{k=1}^{C}
\left(I^o_{y,x,k}-I^c_{y,x,k}\right)^2 .
\end{equation}

Peak signal-to-noise ratio is

\begin{equation}
\operatorname{PSNR}(I^o, I^c) =
10 \log_{10}\left(\frac{1}{\operatorname{MSE}(I^o, I^c)}\right).
\end{equation}

PSNR is sensitive to pixel-level changes and alignment. It is useful for
detecting gross render drift; grounding and navigation behavior require the
separate checks reported later.

\subsection{SSIM}

For local windows with means $\mu_o,\mu_c$, variances $\sigma_o^2,\sigma_c^2$,
and covariance $\sigma_{oc}$, structural similarity is

\begin{equation}
\operatorname{SSIM}(I^o,I^c)=
\frac{(2\mu_o\mu_c+C_1)(2\sigma_{oc}+C_2)}
{(\mu_o^2+\mu_c^2+C_1)(\sigma_o^2+\sigma_c^2+C_2)} .
\end{equation}

The reported SSIM aggregates over windows and pairs. High SSIM means the
matched render pair preserves local luminance and contrast structure under this
metric. Material mechanisms such as clearcoat, procedural texture, and
displacement remain separate diagnostics.

\medskip

\subsection{Feature Similarity}

For CLIP or DINOv2 embeddings $f(I)$, cosine similarity is

\begin{equation}
s_{\cos}(I^o,I^c)=
\frac{f(I^o)^\top f(I^c)}
{\|f(I^o)\|_2\|f(I^c)\|_2}.
\end{equation}

Feature similarity is less tied to exact pixels than PSNR, but it is still a
proxy diagnostic. It can remain high when a localized target cue is lost or
when a coordinate-returning VLM changes point behavior under the same prompt.

\begin{center}
\includegraphics[width=\textwidth,height=0.40\textheight,keepaspectratio]{figures/fig_supplement_grounding_derivation_companion.png}
{\captionsetup{hypcap=false,font=scriptsize,skip=2pt}
\captionof{figure}{Grounding derivation render companion. Target-view
render pairs and the formula-to-render audit strip connect point-in-box and
normalized-coordinate formulas to visible original/noMDL renders. The
coordinate-conversion render wall adds zoom-level target-view anchors for the
same scorer equations, and the target close-up scorer wall shows pass and
diagnostic target views.}}
\label{fig:supp-grounding-derivation-companion}
\end{center}

\subsection{Point-in-Box Grounding}

For each VLM example $j$, let $b_j=(x_{\min},y_{\min},x_{\max},y_{\max})$ be
the target bounding box in image coordinates. Let $p_j=(x_j,y_j)$ be the model
point after parsing and coordinate-frame conversion. A point hit is

\begin{equation}
h_j =
\mathbb{1}\left[
x_{\min}\leq x_j \leq x_{\max}
\land
y_{\min}\leq y_j \leq y_{\max}
\right].
\end{equation}

The hit count is $\sum_j h_j$ over scorable rows. Missing answers, unparseable
coordinates, or answers outside the required protocol are not silently
converted into successful hits.

\subsection{Normalized-1000 Coordinate Contract}

The canonical VLM prompt asks for normalized coordinates in a $1000\times1000$
frame. For an image of width $W$ and height $H$, a parsed normalized point
$(u,v)$ is converted to image pixels by

\begin{equation}
x = \frac{u}{1000}W,\qquad y = \frac{v}{1000}H .
\end{equation}

Raw-image rescoring is retained only as a diagnostic. If a model outputs raw
pixel coordinates despite the normalized prompt, raw scoring can explain model
behavior; the prompt-compliant metric remains the normalized-1000 score.

\subsection{Pair Agreement}

For paired original/converted examples, agreement can be computed on answer
labels or on point-hit statuses. For a binary status $z_j^o,z_j^c\in\{0,1\}$,
pair agreement is

\begin{equation}
a_j = \mathbb{1}[z_j^o=z_j^c],
\qquad
\bar{a}=\frac{1}{M}\sum_{j=1}^{M}a_j .
\end{equation}

Agreement is not equivalent to correctness. Both renders can agree on a miss,
or both can agree on an answer that is categorically correct but
coordinate-inconsistent. This is why the tables report answer, raw point, and
normalized point quantities separately.

\subsection{Bootstrap Paired Intervals}

For a paired metric difference $d_i=m_i^c-m_i^o$, the paired bootstrap samples
indices $i_1,\ldots,i_N$ with replacement and computes

\begin{equation}
\Delta^*=\frac{1}{N}\sum_{r=1}^{N}d_{i_r}.
\end{equation}

The reported confidence interval is the percentile interval over repeated
bootstrap draws. It describes the frozen evidence pool.

\begin{center}
\includegraphics[width=\textwidth,height=0.30\textheight,keepaspectratio]{figures/fig_supplement_navigation_derivation_companion.png}
{\captionsetup{hypcap=false,font=scriptsize,skip=1pt}
\captionof{figure}{Navigation derivation media companion. The formula-to-route
panel links SR, NE, SPL, and TL definitions to metric, route, still, and case
panels.}}
\label{fig:supp-navigation-derivation-companion}
\end{center}

\subsection{Navigation Metrics}

For an episode, success rate is

\begin{equation}
\operatorname{SR}=\mathbb{1}[\operatorname{dist}(p_T,g)\leq \tau],
\end{equation}

where $p_T$ is the final position, $g$ is the goal, and $\tau$ is the official
success threshold. Navigation error is

\begin{equation}
\operatorname{NE}=\operatorname{dist}(p_T,g).
\end{equation}

SPL follows the standard success-weighted path-length form

\begin{equation}
\operatorname{SPL}=S\frac{L^*}{\max(L,L^*)},
\end{equation}

where $S$ is the success indicator, $L$ is the executed path length, and $L^*$
is the reference shortest path length. Trajectory length (TL) records the
executed path length. The main paper reports original/noMDL aggregate means and
paired intervals for the frozen official-scene run.

\medskip

\subsection{Metric Scope Summary}

The derivations above motivate the separate checks. Pixel metrics, feature
metrics, VLM point hits, coordinate compliance, and navigation metrics measure
different objects. A high proxy score can coexist with a prompt-compliance
failure. A successful route can coexist with a material approximation. A
selected visual example can make a mechanism legible without estimating a
population rate.

For this paper, the conversion method is therefore not presented as a
universal semantic compiler. It is presented as a reproducible way to produce
matched original/converted evidence and to expose whether downstream language
grounding or embodied-stack checks still pass in explicit settings.

\begin{center}
\includegraphics[width=\textwidth,height=0.52\textheight,keepaspectratio]{figures/fig_supplement_metric_boundary_bridge.png}
{\captionsetup{hypcap=false}
\captionof{figure}{Metric summary panel. The rows connect each metric
family to visual evidence and to the setting it measures: proxy
metrics inspect appearance, VLM point hits inspect a
prompt-and-box format, material diagnostics inspect covered or selected
mechanisms, and navigation metrics inspect the frozen official-scene run. The
metric-scope lens strip places selected-failure render crops next to a compact
visual summary.}
\label{fig:supp-metric-scope-bridge}}
\end{center}

\clearpage
\twocolumn
